\documentclass[12pt,a4paper]{report}
\usepackage[utf8]{inputenc}
\usepackage[T1]{fontenc}
\usepackage[french]{babel}
\usepackage{graphicx}
\usepackage{geometry}
\usepackage{hyperref}
\usepackage{setspace}
\usepackage{titlesec}
\usepackage{lipsum}

\geometry{margin=2.5cm}
\setstretch{1.2}

\titleformat{\chapter}[hang]{\bfseries\LARGE}{\thechapter.}{1em}{}

\begin{document}

%---------------- PAGE DE TITRE ----------------%
\begin{titlepage}
\centering
\vspace*{3cm}

{\Large Année académique 2025--2026\\[1.5cm]}

{\Huge \textbf{RAPPORT LAB 1}}\\[1cm]

{\large rédigé par :}\\[0.5cm]
{\Large \textbf{NNA Francis Emmanuel Roméo}}\\[0.5cm]
{\large 4\textsuperscript{e} année Cybersécurité et Investigation Numérique}\\[2cm]


\vfill
\textbf{Université  de YAOUNDE I/ École Nationale Superieur Polytechnique Yaounde}\\[0.2cm]
\today

\end{titlepage}

%---------------- TABLE DES MATIÈRES ----------------%
\tableofcontents
\newpage

%---------------- INTRODUCTION ----------------%
\chapter*{Introduction}
\addcontentsline{toc}{chapter}{Introduction}

L'objectif du Lab 1 est de permettre aux étudiants de créer un environnement réseau complexe dans GNS3. 
Ils doivent configurer un réseau sécurisé incluant une machine Windows, une DMZ (zone démilitarisée) 
avec un serveur web sous Linux, un pare-feu et un routeur. 
Ce Lab met en pratique les concepts de segmentation réseau et d'isolation des services critiques. 

Tout au long de ce document, nous allons détailler la mise en œuvre de ce laboratoire jusqu’à son niveau fonctionnel.

%---------------- CONCEPTION ----------------%
\chapter{Conception}

\section{Composition de l’architecture}
Notre infrastructure réseau est composée de :
\begin{itemize}
    \item Un routeur comme équipement de frontière ;
    \item Une DMZ contenant un poste de travail (serveur Ubuntu) qui héberge une application web accessible depuis l’extérieur comme de l’intérieur du réseau d’entreprise ;
    \item Un poste de travail (Windows 10) sur le réseau local, muni d’un antivirus et contenant environ 2 Go de données (fichiers, exécutables, documents, etc.).
\end{itemize}

\section{Schéma de l’architecture}
\begin{figure}[h]
    \centering
    \includegraphics[width=0.85\textwidth]{schema_gns3.png}
    \caption{Shema de L'architecture Reseau}
\end{figure}

\section{Plan d’adressage sur GNS3}

\begin{center}
\begin{tabular}{|l|l|l|l|}
\hline
\textbf{Appareil} & \textbf{Interface} & \textbf{Adresse IP} & \textbf{Masque de sous-réseau} \\ \hline
Ubuntu & ens33 & 192.168.2.31 & 255.255.255.0 \\ \hline
Windows & eth1 & 192.168.1.20 & 255.255.255.0 \\ \hline
Pare-feu (LAN) & em1 & 192.168.1.1 & 255.255.255.0 \\ \hline
Pare-feu (DMZ) & em2 & 192.168.2.30 & 255.255.255.0 \\ \hline
\end{tabular}
\end{center}

%---------------- DÉPLOIEMENT ----------------%
\chapter{Déploiement}

\section{Création et configuration des machines virtuelles}

\subsection{Machine virtuelle Windows 10}
Une machine virtuelle Windows 10 a été créée avec les caractéristiques suivantes :  
\begin{itemize}
    \item Disque dur : 20 Go  
    \item RAM : 2 Go  
\end{itemize}
\section{Machine Virtuelle Windows 10}
\begin{figure}[h]
    \centering
    \includegraphics[width=0.85\textwidth]{windows 10.png}
    \caption{Machine Virtuelle Windows 10}
\end{figure}

Nous avons copié environ 2 Go de données dans les dossiers \textit{Bureau} et \textit{Mes documents}, 
et installé l’antivirus Avast.

\subsection{Machine virtuelle Linux (serveur web)}
Une machine virtuelle Ubuntu a été configurée avec :  
\begin{itemize}
    \item Disque dur : 20 Go  
    \item RAM : 2 Go
\end{itemize}
\section{Machine Virtuelle Windows 10}
\begin{figure}[h]
    \centering
    \includegraphics[width=0.85\textwidth]{ubuntu.png}
    \caption{Machine Virtuelle ubuntu}
\end{figure}
\subsection{Création d’une application web}
Pour simplifier la mise en place, nous avons utilisé le CMS \textbf{WordPress} hébergé sur le serveur Ubuntu.  
Les étapes principales :
\begin{enumerate}
    \item Installation du serveur local XAMPP avec Apache et MySQL.
    \item Extraction de l’archive WordPress dans \texttt{/opt/lampp/htdocs/}.
    \item Lancement du CMS via \url{http://localhost/wordpress}.
\end{enumerate}

\section{Création de l’infrastructure}

\subsection{Installation de GNS3}
L’installation de GNS3 s’est déroulée en deux étapes :
\begin{enumerate}
    \item Installation de la machine virtuelle GNS3.
    \item Installation de l’application GNS3 sur le poste local.
\end{enumerate}

\subsection{Installation et configuration du pare-feu pfSense}
Nous avons utilisé la version 2.7 de pfSense comme pare-feu.  
L’interface LAN (em1) reçoit par défaut l’adresse \texttt{192.168.1.1/24}.  
Depuis la VM Windows, nous nous connectons à l’interface graphique pour configurer les interfaces DMZ et WAN.

\begin{figure}[h]
    \centering
    \includegraphics[width=0.8\textwidth]{pfsense_interface.png}
    \caption{Interface graphique de pfSense après configuration}
\end{figure}

\subsection{Adressage de la machine Windows et test du ping}
L’adresse IP \texttt{192.168.1.20} a été attribuée à la VM Windows.  
Nous ouvrons pfSense via le navigateur avec \url{http://192.168.1.1} 
(identifiants : \texttt{admin / pfsense}).

\subsection{Adressage de la machine Linux et test du ping}
L’utilisateur Ubuntu (mot de passe \texttt{root}) est configuré avec l’adresse \texttt{192.168.2.31}.  
Après lancement de XAMPP, la communication est testée avec un ping vers les autres machines.

\subsection{Test global de l’architecture}
Depuis la VM Windows, nous accédons à \url{http://192.168.2.31/wordpress} et la page d’accueil du site s’affiche correctement.

\begin{figure}[h]
    \centering
    \includegraphics[width=0.85\textwidth]{wordpress_accueil.png}
    \caption{Page d’accueil du site WordPress vue depuis la VM Windows}
\end{figure}

%---------------- CONCLUSION ----------------%
\chapter*{Conclusion}
\addcontentsline{toc}{chapter}{Conclusion}

Nous voici arrivés au terme du Lab 1.  
Le ping entre les différentes machines est fonctionnel et l’application web est accessible depuis la VM Windows.  
Ce projet a permis de mettre en pratique la configuration d’un réseau sécurisé, la segmentation des services 
et le déploiement d’un pare-feu dans un environnement GNS3.

\end{document}
